\section{Uniform local consequences and concrete constructions}\label{dt:sec-8}

The directed criterion has the following simple consequence:

\begin{equation}\label{dt:eq-47}
|n_v|\ge2^{27}\quad\text{at every vertex}
             \quad\Longrightarrow\quad\text{Theorem~\ref{dt:main}}.
\end{equation}

This sufficient constant is independent of the number of vertices,
their degrees and the choice of root. Here is a proof covering all
degrees simultaneously. Choose \(c_v=k_0/2\) at every nonroot vertex
and write \(\delta\) for their common value in \eqref{dt:eq-3}. The elementary
cosine estimate gives

\begin{equation}\label{dt:eq-48}
\chi\ge\frac{r_0^2}{2}-\frac{r_0^4}{24}
       =2^{-9}\frac{3071}{3072}>2^{-10},\qquad
                          \delta>2^{-30}.
\end{equation}

At any vertex, including the root, put
\(\mathcal R_v=\sqrt{d_v^2+n_v^2}\). Then

\begin{equation}\label{dt:eq-49}
A_v\ge\frac{k_0d_v}{2\pi^2},\qquad B_v\le16d_v,
 \qquad C_v\le\frac{\mathcal R_v^2}{8k_0d_v}.
\end{equation}

For the last inequality use \(16\le1/(8k_0)\). Each of the four
branches of the minimum in \eqref{dt:eq-6} consequently has the following lower
bound, respectively:

\begin{equation}\label{dt:eq-50}
\frac\delta2,\qquad
 \frac{t_0k_0d_v}{\pi\mathcal R_v},\qquad
 \frac{k_0d_v}{16\mathcal R_v},\qquad 2^{-12}.
\end{equation}

The second displayed quantity is smaller than the third and fourth.
For the third bound, use
\(\pi|b_v|+t_0|D_v|\le\mathcal R_v(1+t_0/\pi)/2\).
Thus every branch of \eqref{dt:eq-7}, not just its tail branch, is bounded below by

\begin{equation}\label{dt:eq-51}
\eta_v\ge\min\left\{
       \frac{t_0\mathcal R_v\delta}{8\pi},
       \frac{t_0^2 k_0d_v}{4\pi^2}\right\}.
\end{equation}

Indeed the second entry here is smaller than each of the other three
entries in \eqref{dt:eq-7}, by their explicit values and \(t_0\ll r_0\).
On the other hand

\begin{equation}\label{dt:eq-52}
\beta_v=\frac{420\pi}{361}\frac{d_v}{\mathcal R_v^2}
                       <\frac{4d_v}{\mathcal R_v^2}.
\end{equation}

For comparison with half the first entry of \eqref{dt:eq-51}, it suffices that
\(\mathcal R_v^2>64\pi/(t_0\delta)\), since \(d_v\le\mathcal R_v\).
The right side is less than \(2^{52}\), using \(\pi<4\) and \eqref{dt:eq-48}.
For comparison with half the second entry it suffices that
\(\mathcal R_v^2>32\pi^2/(t_0^2k_0)\), whose right side is less
than \(2^{53}\), since \(\pi^2<10\). Both inequalities follow
strictly from \(\mathcal R_v^2\ge n_v^2\ge2^{54}\).
The charge-cut condition in \eqref{dt:eq-8} follows as well.

For the outgoing patch, the sharper elementary estimate is

\begin{equation}\label{dt:eq-53}
\beta_v\le\frac{210\pi}{361|n_v|}
 \le\frac{330}{361}\,2^{-26}
 <\frac{3071}{3072}\,2^{-26}
 \le\frac{k_0}2\chi.
\end{equation}

Here \(\pi<22/7\),
\(d_v/(d_v^2+n_v^2)\le1/(2|n_v|)\), and the displayed rational
comparison prove \eqref{dt:eq-9} for every degree. This completes the corollary
without a numerical scan or a threshold depending on the tree.

The same proof gives a mixed condition, chosen independently at each
vertex:

\begin{equation}\label{dt:eq-53a}
\boxed{\qquad
       d_v\ge2^{28}\quad\text{or}\quad |n_v|\ge2^{27}
               \qquad\text{for every }v.\qquad}
\end{equation}

Under the degree alternative, \(\mathcal R_v^2\ge d_v^2\ge2^{56}\)
proves the same two tail comparisons, while

\begin{equation}\label{dt:eq-53b}
\beta_v\le\frac{420\pi}{361d_v}
             \le\frac{330}{361}\,2^{-26}
                    <\frac{k_0}2\chi
\end{equation}

is the same outgoing patch bound. Neither alternative needs to hold
uniformly over all vertices. In particular, the high-degree vertices
may retain any fixed signed powers, including zero and small nonzero
powers. No sum over the number of such vertices occurs.

Here is a precise extension construction illustrating that freedom.
Start with any fixed ordered circle-tree core whose original endpoint
pairing closes to a knot and has no internal closed component. Retain
all its fixed signed powers. At every core vertex choose one of its
actual product slots next to its crossing box and insert a consecutive
stack of an even number \(B\ge2^{28}\) of new terminal pieces. Every
new leaf has a fixed odd power satisfying

\begin{equation}\label{dt:eq-53c}
|n_{\mathrm{leaf}}|\ge2^{27}+1.
\end{equation}

The physical quarter-turn of an odd two-strand braid has diagonal
pairing \(X=\{02,13\}\); this pairing is invariant under the
quarter-turn. An even ordered stack of these leaves has the through
pairing \(I=\{03,12\}\) and creates no internal circle. Inserting it
in the specified slot therefore preserves the original local endpoint
pairing in that actual frame. This is a pairing calculation, not an
isotopy removing the new tangles. At the original closing box the same
insertion takes place inside its local ball. An added terminal leaf may
then be chosen as the recursive root using the original closed-collar
move with its induced frame reversal and factor order retained.

Thus the resulting diagram is again a knot, while its knot type may
change. Each old vertex now satisfies the degree alternative and each
new leaf satisfies the twist alternative in \eqref{dt:eq-53a}. The resulting actual
knot complement is hyperbolic and satisfies Theorem~\ref{dt:main}. The number and
arrangement of prescribed short-power vertices are arbitrary but fixed;
the number \(B\) does not depend on their total count.

\begin{figure}[t]
\centering\begin{tikzpicture}[x=1cm,y=1cm,>=Stealth,every node/.style={font=\small}]
\node[circle,draw,minimum size=7mm] (v) at (-3,1.5) {$v$};
\node[circle,draw,minimum size=6mm] (l) at (-4,2.7) {$m$};
\node[circle,draw,minimum size=6mm] (r) at (-2,2.7) {$m$};
\draw (v)--node[left] {$T$}(l);
\draw (v)--node[right] {$T$}(r);
\draw[densely dashed] (v)--(-4,1);
\draw[densely dashed] (v)--(-2,1);
\node[align=left,text width=6.7cm] at (2,2)
 {One consecutive pair in the chosen slot.\\Each new power $m$ is odd.\\Repeat the pair $B/2$ times, with $B$ even.};
\begin{scope}[shift={(-3.8,-1.2)}]
\draw[thick] (0,0)--(1,1)--(0,2);
\draw[thick] (1,0)--(0,1)--(1,2);
\draw[dashed,gray] (-.2,1)--(1.2,1);
\foreach \x/\y in {0/0,1/0,0/2,1/2}{\fill (\x,\y) circle(1.5pt);}
\node[left] at (0,0) {$0$};\node[right] at (1,0) {$1$};
\node[right] at (1,2) {$2$};\node[left] at (0,2) {$3$};
\node[left] at (-.3,.5) {$X$};\node[left] at (-.3,1.5) {$X$};
\node at (.5,-.45) {Pairings only};
\end{scope}
\draw[->,thick] (-1.8,-.2)--(-.2,-.2);
\begin{scope}[shift={(.3,-1.2)}]
\draw[thick] (0,0)--(0,2);\draw[thick] (1,0)--(1,2);
\foreach \x/\y in {0/0,1/0,0/2,1/2}{\fill (\x,\y) circle(1.5pt);}
\node[left] at (0,0) {$0$};\node[right] at (1,0) {$1$};
\node[right] at (1,2) {$2$};\node[left] at (0,2) {$3$};
\node at (.5,-.45) {$I=\{03,12\}$};
\end{scope}
\node[align=left,text width=3.2cm] at (4,-.2) {No component closes inside the stack.};
\end{tikzpicture}

\caption{The upper drawing records tangle operations, not a knot
projection. In the construction, each drawn pair represents \(B/2\)
consecutive pairs inserted in one chosen product slot. The lower
pairing calculation is \(X\circ X=I\), with no internal component;
it preserves the original closure's component count, not its knot type.}
\label{dt:leaves-figure}
\end{figure}

The earlier uniform local condition is also a corollary:
\begin{corollary}\label{dt:local}
If \(\max_v d_v/(d_v^2+n_v^2)\le2^{-120}\), then the conclusions
of Theorem~\ref{dt:main} hold.
\end{corollary}
\begin{proof}
If both \(d_v<2^{28}\) and \(|n_v|<2^{27}\), then
\(d_v^2+n_v^2<2^{57}\). Since \(d_v\ge1\), this would give
\(d_v/(d_v^2+n_v^2)>2^{-57}>2^{-120}\). Thus every vertex
satisfies \eqref{dt:eq-53a}. This proof uses the directed theorem,
not a summability hypothesis over vertices.
\end{proof}

\section{A star beyond both previous quantitative tests}\label{dt:sec-9}

Consider an ordered star rooted at its central vertex, of odd degree
\(d=30{,}000{,}001\). Set its power to 2 and every terminal power to the fixed
odd integer \(2^{60}+1\). Choose \(c_v=k_0/2\) on every leaf.

The knot pairing can be checked in the original terminal-root syntax,
even though the finite contraction is evaluated from the center. Label
the boundary points lower left, lower right, upper right, upper left
by \(0,1,2,3\). The even and odd braid pairings are respectively

\begin{equation}\label{dt:eq-54a}
I=\{03,12\},\qquad X=\{02,13\};\qquad
 T(I)=H=\{01,23\},\quad T(X)=X.
\end{equation}

At a terminal root there are \(d-1\) other leaves. Their turned odd
braids have pairing \(X\); their ordered product has pairing \(I\)
because \(d-1\) is even. The central power \(2\) preserves it. Its
actual quarter-turn has pairing \(H\), and stacking the root's odd
braid against \(H\) still gives \(H\), with no internal closed arc.
Cylindrical closure joins \(03\) and \(12\), so it closes \(H\)
into one component. All intermediate products before that last stack
are two through-strand braids. There is no hidden internal component.

The criterion can be verified by rational enclosures without expanding this large
tree. All leaves have the same data. The center has approximately

\begin{equation}\label{dt:eq-54}
\beta_o=1.219\,10^{-7},\qquad
 \eta_o=2.713\,10^{-7},\qquad
       \frac{\beta_o}{\eta_o/2}<0.9.
\end{equation}

Every leaf has a much smaller ratio. The root does not satisfy the
nonroot patch inequality for any \(c>0\), because \(\beta_o>k_0\chi\).
Nevertheless its exact final trace satisfies \eqref{dt:eq-8}. Thus that outgoing
patch must not be charged to a root that sends no message.

Moreover

\begin{equation}\label{dt:eq-55}
\sum_v\frac{d_v}{d_v^2+n_v^2}
       >\frac{30,000,001}{30,000,001^2+4}>2^{-25},
\end{equation}

and its maximum is far above \(2^{-120}\). The example therefore lies beyond
both published quantitative alternatives. Its size and all its powers
are fixed; it is not an N-growing family. The constant sizes are coarse,
and no optimality assertion is intended.

Here are rational enclosures proving the inequalities rather than
inferring them from those decimal values. Put

\begin{equation}\label{dt:eq-56}
\pi_- =3,\quad \pi_+=22/7,\quad c=k_0/2,\quad
 \chi_-=r_0^2/2-r_0^4/24,\quad
 \delta_- =\frac{c\chi_-}{4(1+c)}.
\end{equation}

At the central vertex use the conservative height
\(y_* =\delta_-/2\). It is at most each branch of \eqref{dt:eq-6}. For example,
with

\begin{equation}\label{dt:eq-57}
D_- =\frac d{2\pi_+},\quad D_+=\frac{d+1}{2\pi_-},\quad
 A_- =\frac{dc}{\pi_+^2},\quad
 C_+=d+\frac{\pi_-^{-2}}{4A_-},
\end{equation}

the three non-strip comparisons are

\begin{equation}\label{dt:eq-58}
y_*<\frac{t_0D_-}{4C_+},\qquad
 y_*<\frac{dc}{8(1+t_0D_+)},\qquad y_*^2<\frac c8.
\end{equation}

The strip comparison follows from \(\delta\ge\delta_-\).
Every quantity in \eqref{dt:eq-56}--\eqref{dt:eq-58} is rational. Direct cross-multiplication
also gives

\begin{equation}\label{dt:eq-59}
\frac{2[420\pi_+d/(361(d^2+4))]}
             {t_0D_-y_*/2}<\frac9{10},
 \qquad
 2\frac{420\pi_+d}{361(d^2+4)}
 <\min\left\{\frac{dc}{2},\ A_-r_0^2,
                          \frac{D_-r_0^2}{8}\right\}.
\end{equation}

These compare the upper bound for twice the first-charge drop with
all four lower bounds for \(\eta_o\). The charge cutoff is greater
than two.

For a leaf with \(n=2^{60}+1\), take
\(y_*=t_0k_0/(2n)\). The same direct rational comparisons, using
\(D_-=n/(2\pi_+)\), \(D_+=(n+1)/(2\pi_-)\),
\(A_-=2k_0/\pi_+^2\), and
\(C_+=16+[n/(2\pi_-)]^2/(4A_-)\), put this height below each
branch of \eqref{dt:eq-6}. They give a lower bound greater than \(2^{-49}\)
for every entry of \eqref{dt:eq-7}, while \(\beta_v<2^{-117}\).
The outgoing patch follows from
\(\beta_v<(k_0/2)\chi_-\). This proves the leaf inequalities.
